\documentclass[12pt,a4paper]{article}
\usepackage[utf8]{vietnam}
\usepackage{amsmath,amssymb}
\usepackage{tabularx}
\usepackage{tikz}
\usetikzlibrary{calc,angles,quotes,arrows.meta,patterns}
\usepackage{xcolor}
\usepackage{enumitem}
\usepackage{fancyhdr}
\usepackage{graphicx}
\usepackage[top=1.0cm, bottom=1.0cm, left=1.2cm, right=1.2cm]{geometry}
\usepackage[dethi]{ex_test}
\SAOPTN{widthSA=3.2,heightSA=0.62}

\setlength{\headheight}{15pt}
\setlength{\parskip}{0pt}

\pagestyle{fancy}
\fancyhf{}
\rhead{\textit{Môn Toán 8 -- Đề số 05}}
\lhead{\textit{Bộ đề ôn tập Giữa học kì I}}
\cfoot{\thepage}

\begin{document}

% --- KHUNG TIÊU ĐỀ ĐỀ THI ---
\noindent
\begin{tabularx}{\linewidth}{X r}
\textbf{TRƯỜNG THCS: \dotfill} & \textbf{ĐỀ KIỂM TRA ĐỊNH KỲ GIỮA HỌC KỲ I} \\
\textbf{Họ và tên: \dotfill} & \textbf{NĂM HỌC 2026 -- 2027} \\
\textbf{Lớp: \dotfill \quad SBD: \dotfill} & \textbf{Môn: TOÁN 8} \\
\hline
\end{tabularx}
\vspace{2mm}

\begin{center}
  \textbf{\large ĐỀ ÔN TẬP GIỮA HỌC KỲ I -- ĐỀ SỐ 05}\\
  \textit{Thời gian làm bài: 90 phút (Không kể thời gian phát đề)}
\end{center}
\vspace{1mm}

% --- PHẦN I ---
\noindent\textbf{\large PHẦN I. CÂU TRẮC NGHIỆM NHIỀU PHƯƠNG ÁN LỰA CHỌN (3,0 điểm)}

\textit{Thí sinh trả lời từ câu 1 đến câu 12. Mỗi câu hỏi thí sinh chỉ chọn một phương án duy nhất.}

\vspace{1mm}
\begin{ex}
Trong các biểu thức sau, biểu thức nào không là đa thức?
\choice
{\True $\dfrac{5}{xy}$}
{$12x^2y$}
{$5$}
{$x^2y + 3xy + 1$}
\end{ex}

\begin{ex}
Trong các đơn thức sau, đơn thức nào đồng dạng với đơn thức $-\dfrac{3}{4}x^2y^3$?
\choice
{$-\dfrac{3}{4}x^3y^2$}
{\True $y^3x^2$}
{$2xy$}
{$3x^2y^3z$}
\end{ex}

\begin{ex}
Đơn thức $10x^3y^2z^3$ chia hết cho đơn thức nào sau đây?
\choice
{$5x^3y^3z^2$}
{$-2x^4y^2z$}
{\True $-9x^3yz^2$}
{$2x^3y^3z^3$}
\end{ex}

\begin{ex}
Kết quả phép tính $2x(x - 2y)$ bằng:
\choice
{$2x^2 - 2y^2$}
{$2x^2 - 4y^2$}
{$2x^2 - 2xy$}
{\True $2x^2 - 4xy$}
\end{ex}

\begin{ex}
Kết quả phép tính $(x - 2y)(x + 2y)$ bằng:
\choice
{$x^2 - 2y^2$}
{$(x + 2y)^2$}
{$(x - 2y)^2$}
{\True $x^2 - 4y^2$}
\end{ex}

\begin{ex}
Đa thức $4x^2 - 9$ được phân tích thành nhân tử là:
\choice
{\True $(2x - 3)(2x + 3)$}
{$(4x + 9)(4x - 9)$}
{$(4x - 3)(4x + 3)$}
{$(2x + 9)(2x - 9)$}
\end{ex}

\begin{ex}
Giá trị của biểu thức $x^2 - 2xy + y^2 - 3$ tại $x = -1, y = 2$ là:
\choice
{$-6$}
{\True $6$}
{$12$}
{$-12$}
\end{ex}

\begin{ex}
Tổng số đo các góc trong một tứ giác bằng:
\choice
{$90^\circ$}
{$120^\circ$}
{$180^\circ$}
{\True $360^\circ$}
\end{ex}

\begin{ex}
Tứ giác $ABCD$ là hình bình hành nếu:
\choice
{$\widehat{A} = \widehat{C}$}
{$\widehat{B} = \widehat{D}$}
{$AB \parallel CD$, $AB = AD$}
{\True $\widehat{A} = \widehat{C}$ và $\widehat{B} = \widehat{D}$}
\end{ex}

\begin{ex}
Độ dài cạnh hình thoi có hai đường chéo bằng $6\text{ cm}$ và $8\text{ cm}$ là:
\choice
{$21\text{ cm}$}
{\True $5\text{ cm}$}
{$14\text{ cm}$}
{$\sqrt{11}\text{ cm}$}
\end{ex}

\begin{ex}
Cho hình chóp tam giác đều có diện tích đáy bằng $12\text{ cm}^2$ và chiều cao bằng $5\text{ cm}$. Thể tích của hình chóp đó là:
\choice
{\True $20\text{ cm}^3$}
{$60\text{ cm}^3$}
{$15\text{ cm}^3$}
{$30\text{ cm}^3$}
\end{ex}

\begin{ex}
Kết quả phân tích đa thức $x^3 + 8$ thành nhân tử là:
\choice
{\True $(x + 2)(x^2 - 2x + 4)$}
{$(x + 2)(x^2 + 2x + 4)$}
{$(x - 2)(x^2 + 2x + 4)$}
{$(x + 2)^3$}
\end{ex}


\vspace{1mm}
\noindent\textbf{BẢNG TRẢ LỜI TRẮC NGHIỆM PHẦN I:}
\begin{center}
\begin{tabular}{|c|c|c|c|c|c|c|c|c|c|c|c|c|}
\hline
\textbf{Câu} & \textbf{1} & \textbf{2} & \textbf{3} & \textbf{4} & \textbf{5} & \textbf{6} & \textbf{7} & \textbf{8} & \textbf{9} & \textbf{10} & \textbf{11} & \textbf{12} \\
\hline
\textbf{Chọn} & \rule{0pt}{0.65cm} & & & & & & & & & & & \\
\hline
\end{tabular}
\end{center}
\vspace{2mm}

\setcounter{ex}{0}
\noindent\textbf{\large PHẦN II. CÂU TRẮC NGHIỆM ĐÚNG SAI (2,0 điểm)}

\textit{Thí sinh trả lời câu 1 và câu 2. Trong mỗi ý a), b), c), d) ở mỗi câu, thí sinh chọn đúng hoặc sai.}

\vspace{1mm}
\begin{ex}
Cho biểu thức $P = \left(\dfrac{1}{x - 2} - \dfrac{1}{x + 2}\right) \cdot \dfrac{x^2 - 4}{4}$ với $x \neq \pm 2$. Xét tính đúng sai của các mệnh đề sau:
\choiceTF
{\True Mẫu thức chung của hai phân thức trong ngoặc là $x^2 - 4$}
{\True Kết quả rút gọn biểu thức trong ngoặc là $\dfrac{4}{x^2 - 4}$}
{\True Giá trị của biểu thức $P$ sau khi rút gọn bằng $1$}
{Tại $x = 3$, giá trị của biểu thức $P$ bằng $3$}
\end{ex}

\begin{ex}
Cho phương trình $(x - 3)(x + 3) - x(x - 2) = 5$. Xét tính đúng sai của các mệnh đề sau:
\choiceTF
{\True Khai triển tích $(x - 3)(x + 3)$ ta được $x^2 - 9$}
{\True Biến đổi vế trái của phương trình ta được $2x - 9$}
{\True Phương trình có nghiệm duy nhất $x = 7$}
{Nghiệm $x = 7$ không là nghiệm của phương trình $3x - 5 = 16$}
\end{ex}


\noindent\textbf{BẢNG TRẢ LỜI PHẦN II:}
\begin{center}
\begin{tabular}{|c|c|c|c|c||c|c|c|c|c|}
\hline
\textbf{Câu 1} & \textbf{a)} & \textbf{b)} & \textbf{c)} & \textbf{d)} & \textbf{Câu 2} & \textbf{a)} & \textbf{b)} & \textbf{c)} & \textbf{d)} \\
\hline
\textbf{Đ/S} & \rule{0pt}{0.6cm} & & & & \textbf{Đ/S} & & & & \\
\hline
\end{tabular}
\end{center}
\vspace{2mm}

\setcounter{ex}{0}
\noindent\textbf{\large PHẦN III. CÂU HỎI TRẢ LỜI NGẮN (2,0 điểm)}

\textit{Thí sinh trả lời câu 1 và câu 2. Điền kết quả vào ô trống tương ứng.}

\vspace{1mm}
\begin{ex}
Tìm $x$, biết: $(x + 2)^2 - x(x - 3) = 14$. \shortans{$x = \dfrac{10}{7}$}
\end{ex}

\begin{ex}
Cho hình thang cân $ABCD$ có $\widehat{A} = 105^\circ$. Tính số đo góc $\widehat{C}$? \shortans{$75^\circ$}
\end{ex}


\setcounter{ex}{0}
\noindent\textbf{\large PHẦN IV. TỰ LUẬN (3,0 điểm)}

\textit{Thí sinh trình bày chi tiết lời giải vào giấy làm bài.}

\vspace{1mm}
\begin{ex} (1,0 điểm)
Thực hiện phép tính (rút gọn nếu có thể):
\begin{enumerate}[label=\textbf{\alph*)}, leftmargin=0.8cm, itemsep=2pt]
\item (0,35 điểm) $(2x^2 + 3xy) + (x^2 - xy)$
\item (0,35 điểm) $(2x + y)(x - 3y)$
\item (0,3 điểm) $(9x^2y^3 - 6x^3y^2 + 4x^2y^2) : (3x^2y)$
\end{enumerate}
\end{ex}
\vspace{0.5mm}

\begin{ex} (0,75 điểm)
Cho hai đa thức: $A = x^7 - 4x^3y^2 - 5xy + 7$ và $B = x^7 + 5x^3y^2 - 3xy - 3$.
\begin{enumerate}[label=\textbf{\alph*)}, leftmargin=0.8cm, itemsep=2pt]
\item (0,4 điểm) Tìm đa thức $C$ sao cho $C = A + B$.
\item (0,35 điểm) Tìm đa thức $D$ sao cho $A + D = B$.
\end{enumerate}
\end{ex}
\vspace{0.5mm}

\begin{ex} (0,75 điểm)
Cho tam giác $ABC$ vuông cân tại $A$, đường trung tuyến $AM$. Gọi $N$ là trung điểm của $AC$. Vẽ điểm $D$ đối xứng với điểm $M$ qua $N$.
\begin{enumerate}[label=\textbf{\alph*)}, leftmargin=0.8cm, itemsep=2pt]
\item (0,35 điểm) Chứng minh $AM = MC$.
\item (0,4 điểm) Chứng minh tứ giác $ADCM$ là hình vuông.
\end{enumerate}
\end{ex}
\vspace{0.5mm}

\begin{ex} (0,5 điểm)
Tìm giá trị nhỏ nhất của biểu thức: $M = 9x^2 + 6y^2 + 18x - 12xy - 12y - 27$.
\begin{enumerate}[label=\textbf{\alph*)}, leftmargin=0.8cm, itemsep=2pt]
\item (0,5 điểm) Tìm giá trị nhỏ nhất của $M$.
\end{enumerate}
\end{ex}
\vspace{0.5mm}


\begin{center}
\rule{4cm}{0.4pt}\\
\textit{--- HẾT ---}
\end{center}
\end{document}
